\documentclass{lecture-kuznia}
\usepackage{comment}

\def\data{01.10.2026}
\def\place{Poręba Wielka}
\def\title{Zaawansowane grafy}
\def\group{Finaliści}
\def\author{Autor: Stefan Świerczewski}
\def\lecturer{Prowadzący: Stefan Świerczewski}

\ifshowsolutions
\ExplSyntaxOn
\tl_new:N \g_trudne_grafy_rozwiazania_tl
\newcounter{rozwiazanie}
\NewDocumentEnvironment{solutionbox}{+b}{
	\tl_gput_right:Nn \g_trudne_grafy_rozwiazania_tl {
		\par\medskip
		\stepcounter{rozwiazanie}
		\begingroup\color{black!85}
		\emergencystretch=1em
		\noindent\textbf{Rozwiązanie~\therozwiazanie.}\quad\allowbreak #1
		\hfill$\blacksquare$\par\medskip
		\endgroup
	}
}{}
\NewDocumentCommand{\drukujrozwiazania}{}{
	\tl_use:N \g_trudne_grafy_rozwiazania_tl
}
\ExplSyntaxOff
\else
\excludecomment{solutionbox}
\newcommand{\drukujrozwiazania}{}
\fi

\newenvironment{theoryproof}{\begin{proof}}{\end{proof}}

\newcounter{zadanie}
\newenvironment{zadanie}[1]{%
	\refstepcounter{zadanie}%
	\par\medskip\noindent
	\textbf{Zadanie \thezadanie.}\quad
}{\par\medskip}

\newcommand{\degG}[1]{\deg(#1)}
\newcommand{\AoPSGraphs}{\href{https://services.artofproblemsolving.com/download.php?id=YXR0YWNobWVudHMvZS9lL2JhMjY3OGNlODZjMTYyZmZjOWMwNTdiMWQ0MTIyNzFiZGEzMzJkLnBkZg==\&rn=T2x5bXBpYWRDb21iaW5hdG9yaWNzQ2hhcHRlcjgucGRm}{AoPS: Pranav Sriram, \textit{Olympiad Combinatorics --- Graph Theory}}}

\begin{document}

\begin{center}
  {\Huge\textbf{\title}}

  \medskip
  {\large 11 problemów o ekstremum, przeszukiwaniu i nierównościach}
\end{center}

\fontsize{12}{16}\selectfont

\section{Narzędzia}

\subsection*{Oznaczenia}

\begingroup
\renewcommand{\arraystretch}{1.18}
\begin{tabularx}{\textwidth}{@{}lX@{}}
	$V(G)$, $E(G)$
	& odpowiednio zbiór wierzchołków i zbiór krawędzi grafu $G$;
	piszemy $n=|V(G)|$ oraz $m=|E(G)|$;\\
	$d_G(v)$ lub $d(v)$
	& stopień wierzchołka $v$, czyli liczba jego sąsiadów;\\
	$d^+(v)$, $d^-(v)$
	& stopień wychodzący i stopień wchodzący w grafie skierowanym;\\
	$\delta(G)$, $\Delta(G)$
	& odpowiednio najmniejszy i największy stopień wierzchołka w $G$;\\
	$\operatorname{avgdeg}(G)=2m/n$
	& średni stopień w grafie; zawsze $\delta(G)\leqslant\operatorname{avgdeg}(G)
		\leqslant\Delta(G)$;\\
	$N(v)$, $N[v]$
	& sąsiedztwo otwarte wierzchołka $v$ oraz sąsiedztwo domknięte
	$N(v)\cup\{v\}$;\\
	$N(X)$
	& zbiór wierzchołków spoza $X$, które mają sąsiada w $X$;\\
	$d(u,v)$, $S_i(v)$
	& odległość między $u$ i $v$ oraz warstwa
	$S_i(v)=\{x:d(v,x)=i\}$;\\
	$G[X]$, $G-X$
	& podgraf indukowany przez $X$ oraz graf po usunięciu wierzchołków $X$;\\
	$\overline G$
	& dopełnienie grafu $G$;\\
	$\alpha(G)$, $\omega(G)$
	& największy rozmiar odpowiednio zbioru niezależnego i kliki;\\
	$\nu(G)$, $\tau(G)$
	& największy rozmiar skojarzenia oraz najmniejszy rozmiar pokrycia
	wierzchołkowego;\\
	$\chi(G)$, $\operatorname{girth}(G)$
	& liczba chromatyczna oraz długość najkrótszego cyklu;\\
	$t_k(G)$
	& liczba klik $K_k$ zawartych w $G$.
\end{tabularx}
\endgroup

Jeżeli graf jest ustalony, indeks $G$ zwykle pomijamy. Symbol $d$ bez
argumentu może też oznaczać pomocniczą liczbę całkowitą; z kontekstu będzie
zawsze jasne, czy chodzi o parametr, czy o stopień $d(v)$.

\subsection{Gęstość, rdzeń i długie ścieżki}

\begin{lemma}[rdzeń o dużym stopniu]
	Jeżeli $m\geqslant dn$, to $G$ zawiera niepusty podgraf $H$ o
	$\delta(H)\geqslant d$.
\end{lemma}

\begin{theoryproof}
	Usuwamy wierzchołki stopnia mniejszego niż $d$. Gdyby proces usunął cały
	graf, każdą krawędź policzylibyśmy przy usunięciu jej pierwszego końca, a
	łącznie usunęlibyśmy mniej niż $dn$ krawędzi.
\end{theoryproof}

\begin{theorem}[Erdős--Gallai]
	Jeżeli graf na $n$ wierzchołkach nie zawiera ścieżki złożonej z $k$
	krawędzi, to
	\[
		m\leqslant\frac{k-1}{2}n.
	\]
	W szczególności średni stopień większy niż $k-1$ wymusza ścieżkę długości
	co najmniej $k$.
\end{theorem}

\begin{theoryproof}
	Najpierw zauważmy, że każdy spójny graf o minimalnym stopniu $\delta$ ma
	ścieżkę o co najmniej $\min\{n,2\delta+1\}$ wierzchołkach. Istotnie, niech
	$P=v_1\ldots v_s$ będzie najdłuższą ścieżką. Jeżeli $s<n$, to nie istnieje
	indeks $i$ taki, że jednocześnie
	\[
		v_1v_{i+1}\in E(G)
		\qquad\text{i}\qquad
		v_iv_s\in E(G).
	\]
	W przeciwnym razie wierzchołki $P$ tworzyłyby cykl, który dzięki spójności
	grafu i istnieniu wierzchołka poza $P$ można byłoby otworzyć do dłuższej
	ścieżki. Zbiory indeksów odpowiadające sąsiadom $v_1$ i $v_s$ są więc
	rozłączne, skąd $s-1\geqslant d(v_1)+d(v_s)\geqslant2\delta$.

	Dowodzimy teraz twierdzenia indukcyjnie po $n$. Dla grafu niespójnego
	stosujemy założenie indukcyjne do każdej składowej i sumujemy otrzymane
	oszacowania, możemy więc założyć spójność. Dla $n\leqslant k$ teza wynika z
	$m\leqslant\binom n2\leqslant(k-1)n/2$. Niech zatem $n\geqslant k+1$.
	Gdyby $\delta(G)>(k-1)/2$, powyższa obserwacja dałaby ścieżkę o co najmniej
	$k+1$ wierzchołkach. Istnieje więc wierzchołek $v$ o
	$d(v)\leqslant(k-1)/2$. Z założenia indukcyjnego
	\[
		m-d(v)\leqslant\frac{k-1}{2}(n-1),
	\]
	a po dodaniu $d(v)$ otrzymujemy żądane $m\leqslant(k-1)n/2$.
\end{theoryproof}

\subsection{Dwa oszacowania ekstremalne}

Niech $T_r(n)$ oznacza pełny graf $r$-dzielny, którego części różnią się
rozmiarami o co najwyżej $1$.

\begin{theorem}[Turán]
	Jeżeli $G$ nie zawiera kliki $K_{r+1}$, to
	\[
		m\leqslant e(T_r(n))
		\leqslant\left(1-\frac1r\right)\frac{n^2}{2}.
	\]
	Pierwsza granica jest ostra i równość zachodzi dla grafu Turána $T_r(n)$.
\end{theorem}

\begin{theoryproof}
	Oznaczmy $t_r(n)=e(T_r(n))$ i przeprowadźmy indukcję po $r+n$. Dla
	$r\geqslant n$ teza jest oczywista, bo $t_r(n)=\binom n2$, a dla $r=1$
	graf nie ma krawędzi. Zauważmy też, że spośród pełnych grafów $r$-dzielnych
	najwięcej krawędzi ma ten o możliwie równych częściach: przeniesienie
	wierzchołka z części większej do mniejszej zmniejsza liczbę par leżących w
	jednej części. W szczególności $t_{r-1}(n)\leqslant t_r(n)$.

	Jeżeli $G$ nie zawiera $K_r$, to z założenia indukcyjnego
	$m\leqslant t_{r-1}(n)\leqslant t_r(n)$. W przeciwnym razie wybieramy klikę
	$C\simeq K_r$. Każdy wierzchołek spoza $C$ ma w $C$ co najwyżej $r-1$
	sąsiadów, gdyż $G$ nie zawiera $K_{r+1}$. Dlatego liczba krawędzi mających
	co najmniej jeden koniec w $C$ nie przekracza
	\[
		\binom r2+(n-r)(r-1).
	\]
	Graf $G-C$ nadal nie zawiera $K_{r+1}$, więc indukcja daje
	\[
		m\leqslant t_r(n-r)+\binom r2+(n-r)(r-1)=t_r(n).
	\]
	Ostatnia równość wynika z usunięcia po jednym wierzchołku z każdej części
	$T_r(n)$. Sam graf $T_r(n)$ nie zawiera $K_{r+1}$, zatem oszacowanie jest
	ostre. Nierówność
	$t_r(n)\leqslant(1-1/r)n^2/2$ wynika z
	\[
		t_r(n)=\frac12\left(n^2-\sum_{i=1}^r n_i^2\right)
		\leqslant\frac12\left(n^2-\frac{n^2}{r}\right),
	\]
	gdzie $n_1+\dots+n_r=n$ są rozmiarami części.
\end{theoryproof}

Drugie oszacowanie zamienia lokalne stopnie na duży zbiór niezależny.

\begin{theorem}[Caro--Wei]
	Każdy graf $G$ spełnia
	\[
		\alpha(G)\geqslant
		\sum_{v\in V(G)}\frac1{d(v)+1}
		\geqslant\frac{n^2}{2m+n}.
	\]
\end{theorem}

\begin{theoryproof}
	Ustawiamy wierzchołki losowo i wybieramy $v$, gdy stoi przed wszystkimi
	swoimi sąsiadami. Wybrane wierzchołki są niezależne, a $v$ zostaje wybrany
	z prawdopodobieństwem $1/(d(v)+1)$. Pierwsza nierówność wynika z liniowości
	wartości oczekiwanej, druga zaś z nierówności Cauchy'ego.
\end{theoryproof}

\subsection{Skojarzenia i dualność}

\begin{theorem}[Hall]
	Graf dwudzielny o częściach $A,B$ ma skojarzenie pokrywające $A$ wtedy i
	tylko wtedy, gdy
	\[
		|N(X)|\geqslant|X|\qquad\text{dla każdego }X\subseteq A.
	\]
\end{theorem}

\begin{theoryproof}
	Konieczność jest natychmiastowa: różne wierzchołki $X\subseteq A$ są
	skojarzone z różnymi wierzchołkami $N(X)$.

	Dla dostateczności prowadzimy indukcję po $|A|$. Jeżeli dla każdego
	niepustego, właściwego $X\subset A$ zachodzi mocniejsza nierówność
	$|N(X)|\geqslant|X|+1$, wybieramy dowolną krawędź $ab$, usuwamy $a$ i $b$
	i stosujemy indukcję. Usunięcie $b$ zmniejsza każde sąsiedztwo o co najwyżej
	$1$, więc warunek Halla pozostaje spełniony.

	W przeciwnym razie istnieje niepusty, właściwy $X\subset A$ taki, że
	$|N(X)|=|X|$. Indukcja daje skojarzenie $X$ z całym $N(X)$. W grafie po
	usunięciu $X\cup N(X)$ także zachodzi warunek Halla: gdyby dla pewnego
	$Y\subseteq A\setminus X$ jego sąsiedztwo poza $N(X)$ miało mniej niż
	$|Y|$ elementów, to
	\[
		|N(X\cup Y)|<|X|+|Y|,
	\]
	co przeczyłoby założeniu. Drugie skojarzenie, otrzymane indukcyjnie, pokrywa
	$A\setminus X$. Suma obu skojarzeń pokrywa całe $A$.
\end{theoryproof}

\begin{theorem}[Kőnig]
	W grafie dwudzielnym maksymalny rozmiar skojarzenia $\nu(G)$ jest równy
	minimalnemu rozmiarowi pokrycia wierzchołkowego. Równoważnie,
	\[
		\alpha(G)+\nu(G)=n.
	\]
\end{theorem}

\begin{theoryproof}
	Niech $A,B$ będą częściami grafu, a $M$ --- skojarzeniem maksymalnego
	rozmiaru. Zaczynamy we wszystkich nieskojarzonych wierzchołkach $A$ i
	idziemy ścieżkami naprzemiennymi: z $A$ do $B$ krawędziami spoza $M$, a z
	$B$ do $A$ krawędziami należącymi do $M$. Niech $Z$ będzie zbiorem
	osiągniętych wierzchołków.

	Zbiór
	\[
		C=(A\setminus Z)\cup(B\cap Z)
	\]
	jest pokryciem wierzchołkowym. Gdyby krawędź $ab$, gdzie $a\in A$ i
	$b\in B$, nie miała końca w $C$, mielibyśmy $a\in Z$ i $b\notin Z$.
	Krawędź spoza $M$ pozwoliłaby dojść z $a$ do $b$, a jeżeli $ab\in M$, to
	$b$ musiał już wystąpić na ścieżce naprzemiennej prowadzącej do $a$.
	W obu przypadkach otrzymujemy sprzeczność.

	Żaden osiągnięty wierzchołek $B$ nie jest wolny, bo wtedy znaleźlibyśmy
	ścieżkę powiększającą $M$. Każda krawędź $M$ ma dokładnie jeden koniec w
	$C$: koniec w $B$, gdy oba końce należą do $Z$, oraz koniec w $A$, gdy żaden
	nie należy do $Z$. Stąd $|C|=|M|$. Każde pokrycie musi zawierać co najmniej
	jeden koniec każdej krawędzi skojarzenia, więc ma rozmiar co najmniej
	$|M|$. Otrzymujemy $\tau(G)=\nu(G)$. Ponieważ dopełnienie pokrycia
	wierzchołkowego jest zbiorem niezależnym i odwrotnie,
	$\alpha(G)+\tau(G)=n$, co daje również $\alpha(G)+\nu(G)=n$.
\end{theoryproof}

\subsection{Planarność: wzór Eulera i nierówności}

Jeżeli spójny graf płaski ma $n$ wierzchołków, $m$ krawędzi i $f$ ścian, to
\[
	n-m+f=2.
\]
Jeżeli każda ściana ma brzeg długości co najmniej $q$, to z podwójnego
zliczenia krawędzi na brzegach ścian otrzymujemy
\[
	qf\leqslant2m,
	\qquad
	m\leqslant\frac{q}{q-2}(n-2).
\]
W szczególności dla prostego grafu planarnego o $n\geqslant3$ mamy
\[
	m\leqslant3n-6,
\]
a jeżeli graf jest dodatkowo dwudzielny, to
\[
	m\leqslant2n-4.
\]
Pierwsza z tych nierówności razem z $\sum_vd(v)=2m$ daje przydatny zapis
\[
	\sum_{v\in V(G)}(6-d(v))=6n-2m\geqslant12.
\]

\begin{theoryproof}
	Najpierw liczymy na dwa sposoby pary $(F,e)$, w których $e$ jest bokiem
	brzegu ściany $F$. Każda krawędź ma dwie strony, więc pojawia się w tym
	zliczeniu dwa razy; dla mostu obie strony należą do tej samej ściany. Zatem
	\[
		\sum_F\ell(F)=2m.
	\]
	Jeżeli każdy brzeg ma długość co najmniej $q$, to
	\[
		qf\leqslant\sum_F\ell(F)=2m.
	\]
	Ze wzoru Eulera mamy $f=2-n+m$. Po podstawieniu
	\[
		q(2-n+m)\leqslant2m,
	\]
	czyli
	\[
		(q-2)m\leqslant q(n-2).
	\]
	Stąd $m\leqslant q(n-2)/(q-2)$. Dla grafu prostego przyjmujemy $q=3$, a
	dla grafu prostego i dwudzielnego $q=4$.

	Podobnie podwójnie liczymy pary $(v,e)$, w których $v$ jest końcem krawędzi
	$e$. Każda krawędź ma dwa końce, więc
	\[
		\sum_vd(v)=2m.
	\]
	Z nierówności $m\leqslant3n-6$ otrzymujemy
	\[
		\sum_v(6-d(v))=6n-2m\geqslant12.
	\]
\end{theoryproof}

\subsection{Lokalizacja klik}

Przy zliczaniu klik warto przechodzić do grafów indukowanych przez
sąsiedztwa. Jeśli $t_k(G)$ oznacza liczbę klik $K_k$ w $G$, to
\[
	\sum_{v\in V(G)}t_{k-1}(G[N(v)])=k\,t_k(G).
\]
Tożsamość zachowuje dokładne stałe i dobrze współpracuje z Cauchym oraz z
oszacowaniami ekstremalnymi.

\begin{theoryproof}
	Liczymy na dwa sposoby pary $(v,Q)$, gdzie $Q$ jest kliką $K_k$, a
	$v\in V(Q)$ jest jednym z jej wyróżnionych wierzchołków. Każda klika $Q$
	daje dokładnie $k$ par, więc liczba par wynosi $k\,t_k(G)$.

	Z drugiej strony ustalmy $v$. Po usunięciu $v$ z kliki $Q$ pozostaje klika
	$K_{k-1}$ złożona z sąsiadów $v$, czyli zawarta w $G[N(v)]$. Każda taka
	klika wraz z $v$ daje dokładnie jedną klikę $K_k$. Dla ustalonego $v$ mamy
	więc $t_{k-1}(G[N(v)])$ par. Po zsumowaniu po wszystkich $v$ otrzymujemy
	lewą stronę dowodzonej tożsamości.
\end{theoryproof}

\newpage
\section{Zadania}

\begin{zadanie}{łatwe} % reformulacja ELMO Shortlist 2011 C7
	Niech $t\geqslant1$. Udowodnij, że każdy graf $G$ spełniający
	\[
		\operatorname{avgdeg}(G)\geqslant2(t-1)
	\]
	zawiera jako podgraf każde drzewo o $t$ wierzchołkach.
\end{zadanie}

\begin{solutionbox}
	Niech $n=|V(G)|$ i $m=|E(G)|$. Z założenia
	\[
		m=\frac{n\operatorname{avgdeg}(G)}2\geqslant(t-1)n.
	\]
	Lemat o rdzeniu daje niepusty podgraf $H\subseteq G$ o
	$\delta(H)\geqslant t-1$.

	Pokażemy indukcyjnie po $t$, że taki graf $H$ zawiera każde drzewo $T$ na
	$t$ wierzchołkach. Dla $t=1$ jest to oczywiste. Dla $t\geqslant2$ usuwamy
	z $T$ liść $x$ i oznaczamy jego sąsiada przez $y$. Z założenia indukcyjnego
	osadzamy $T-x$ w $H$. Obraz wierzchołka $y$ ma co najmniej $t-1$ sąsiadów,
	a spośród nich najwyżej $t-2$ zostało już użytych. Istnieje więc wolny
	sąsiad, w którym umieszczamy $x$. Otrzymujemy kopię całego $T$ w $G$.
\end{solutionbox}

\begin{zadanie}{rozgrzewka} % turnieje; dwa zastosowania ekstremum
	Turniej to orientacja grafu pełnego. Udowodnij, że:
	\begin{enumerate}
		\item w każdym turnieju istnieje wierzchołek, z którego do każdego innego
		można dojść ścieżką skierowaną długości co najwyżej $2$;
		\item jeżeli turniej zawiera cykl skierowany, to zawiera skierowany
		trójkąt.
	\end{enumerate}
\end{zadanie}

\begin{solutionbox}
	W pierwszej części wybierzmy wierzchołek $v$ o największym stopniu
	wyjściowym. Jeśli $v$ nie prowadzi bezpośrednio do $u$, to $u\to v$.
	Gdyby nie istniał wierzchołek $w$ taki, że $v\to w\to u$, wierzchołek $u$
	prowadziłby do $v$ i do każdego wychodzącego sąsiada $v$. Miałby więc
	stopień wyjściowy większy od stopnia $v$, sprzecznie z wyborem $v$.

	W drugiej części wybierzmy najkrótszy cykl skierowany
	$v_1\to v_2\to\dots\to v_m\to v_1$. Jeśli $m>3$, rozpatrzmy krawędź
	między $v_1$ i $v_3$. Kierunek $v_3\to v_1$ daje skierowany trójkąt, a
	kierunek $v_1\to v_3$ daje krótszy cykl. Obie możliwości przeczą
	minimalności $m$, zatem $m=3$.
\end{solutionbox}




\begin{zadanie}{łatwe} % Vietnam TST 2001
	Klub ma $42$ członków. W każdym zbiorze $31$ członków istnieją chłopak i
	dziewczyna, którzy się znają. Udowodnij, że można utworzyć $12$ rozłącznych
	par chłopak--dziewczyna, w których osoby się znają.
\end{zadanie}

\begin{solutionbox}
	Tworzymy graf dwudzielny: po lewej są chłopcy, po prawej dziewczęta, a
	krawędź oznacza znajomość. Warunek mówi, że każdy zbiór niezależny ma co
	najwyżej $30$ wierzchołków, czyli $\alpha(G)\leqslant30$.

	Z twierdzenia Kőniga
	\[
		\nu(G)=|V(G)|-\alpha(G)\geqslant42-30=12.
	\]
	Istnieje więc skojarzenie złożone z co najmniej $12$ krawędzi, a wybrane
	krawędzie są żądanymi parami.
\end{solutionbox}

\begin{zadanie}{łatwe} % MOP 2008
	Niech $n\geqslant3$. Krawędzie grafu pełnego $K_n$ pokolorowano tak, że żaden kolor nie został
	użyty na więcej niż $n-2$ krawędziach. Udowodnij, że istnieje trójkąt,
	którego trzy krawędzie mają trzy różne kolory.
\end{zadanie}

\begin{solutionbox}
	Załóżmy przeciwnie. Wybierzmy największą spójną składową utworzoną przez
	krawędzie jednego koloru; nazwijmy ją $X$, a jej kolor czerwonym.

	Gdyby $v\notin X$, to dla dowolnej czerwonej krawędzi $xy$ w $X$ krawędzie
	$vx$ i $vy$ nie mogą być czerwone. Nie mogą też mieć różnych kolorów, bo
	$vxy$ byłby tęczowym trójkątem. Mają więc ten sam kolor. Idąc po czerwonych
	ścieżkach w $X$, widzimy, że wszystkie krawędzie z $v$ do $X$ mają jeden
	kolor. Wtedy $X\cup\{v\}$ leży w większej jednokolorowej składowej, co
	przeczy wyborowi $X$. Zatem $X$ zawiera wszystkie $n$ wierzchołków.

	Spójny graf na $n$ wierzchołkach ma co najmniej $n-1$ krawędzi. Kolor
	czerwony wystąpił więc co najmniej $n-1$ razy, wbrew założeniu.
\end{solutionbox}

\begin{zadanie}{średnie} % krzywizna Eulera; trening finałowy
	Niech $G$ będzie prostym, spójnym grafem planarnym o minimalnym stopniu co
	najmniej $5$. Udowodnij, że $G$ ma co najmniej $12$ wierzchołków stopnia
	dokładnie $5$. Wykaż ponadto, że jeżeli takich wierzchołków jest dokładnie
	$12$, to każda ściana dowolnego rysunku płaskiego $G$ jest trójkątem, a
	każdy pozostały wierzchołek ma stopień $6$.
\end{zadanie}

\begin{solutionbox}
	Oznaczmy przez $n,m,f$ liczby wierzchołków, krawędzi i ścian $G$. Ze wzoru
	Eulera i nierówności $3f\leqslant2m$ mamy $m\leqslant3n-6$. Zatem
	\[
		\sum_{v\in V(G)}(6-d(v))=6n-2m\geqslant12. \tag{1}
	\]

	Niech $x$ oznacza liczbę wierzchołków stopnia $5$. Ponieważ wszystkie
	stopnie są co najmniej $5$, lewa strona (1) jest równa
	\[
		x-\sum_{d(v)\geqslant7}(d(v)-6)\leqslant x.
	\]
	Stąd $x\geqslant12$.

	Jeżeli $x=12$, wszystkie nierówności muszą być równościami. Nie ma więc
	wierzchołka stopnia większego niż $6$, a każdy wierzchołek inny niż
	wyróżnione dwanaście ma stopień $6$. Ponadto $m=3n-6$. Ze wzoru Eulera
	$f=2n-4$, czyli $3f=2m$. Ponieważ każda ściana ma długość co najmniej $3$,
	równość ta oznacza, że wszystkie ściany są trójkątami.
\end{solutionbox}

\begin{zadanie}{średnie} % Croatian TST 2011
	Na przyjęciu jest $n$ osób. Wśród dowolnych czterech osób istnieją albo
	trzy parami zaprzyjaźnione, albo trzy parami nieznające się. Udowodnij, że
	uczestników można rozdzielić na zbiory $A,B$ tak, aby $A$ był kliką, a $B$
	zbiorem niezależnym.
\end{zadanie}

\begin{solutionbox}
	Niech $A$ będzie największą kliką, a $B=V(G)\setminus A$. Wystarczy dowieść,
	że $B$ jest niezależny. Przypuśćmy więc, że $v_1v_2$ jest krawędzią w $B$.
	Ponieważ $A$ jest maksymalna, istnieje $u_1\in A$ niepołączony z $v_1$.

	Jeśli $u_1$ nie jest połączony także z $v_2$, to dla każdego
	$u\in A\setminus\{u_1\}$ w czwórce $u,u_1,v_1,v_2$ nie może być niezależnej
	trójki: są w niej dwie rozłączne krawędzie $uu_1$ i $v_1v_2$. Musi więc być
	trójkąt, a jedyną możliwością jest $uv_1v_2$. Wtedy
	$(A\setminus\{u_1\})\cup\{v_1,v_2\}$ jest większą kliką.

	Jeśli $u_1v_2$ jest krawędzią, to dla każdego $u\in A\setminus\{u_1\}$ ta sama czwórka nie
	zawiera niezależnej trójki, a każdy możliwy trójkąt wymusza krawędź
	$uv_2$. Zatem $v_2$ łączy się ze wszystkimi w $A$, więc $A\cup\{v_2\}$ jest
	większą kliką. Oba przypadki dają sprzeczność.
\end{solutionbox}

\begin{zadanie}{średnie} % AoPS handout, ćw. 29; uogólnienie zadania Rosja 2001
	Drzewo ma dokładnie $2n$ liści, gdzie $n\geqslant2$. Dodaj do niego $n$
	krawędzi tak, aby po usunięciu dowolnej krawędzi otrzymany graf nadal był
	spójny.
\end{zadanie}

\begin{solutionbox}
	Ukorzeńmy drzewo w wierzchołku stopnia co najmniej $2$ i wypiszmy liście
	$\ell_1,\ell_2,\dots,\ell_{2n}$ w kolejności DFS. Dodajmy krawędzie
	\[
		\ell_i\ell_{i+n}\qquad(1\leqslant i\leqslant n).
	\]
	Usunięcie jednej z nowych krawędzi nie rozspaja grafu, bo pozostaje całe
	początkowe drzewo.

	Rozważmy starą krawędź $e$. Po jej usunięciu jedna ze składowych jest
	poddrzewem ukorzenionym, więc jej liście tworzą niepusty przedział $I$ w
	cyklicznym porządku DFS. Jest to przedział właściwy, ponieważ po obu stronach
	$e$ leży liść początkowego drzewa. Pokażemy, że któraś z dodanych krawędzi
	przecina ten podział.

	Jeżeli $|I|\leqslant n$, wybierzmy liść $\ell_i\in I$. Liść położony
	naprzeciwko niego, czyli przesunięty cyklicznie o $n$ pozycji, nie należy do
	$I$ i jest połączony z $\ell_i$. Jeżeli $|I|>n$, stosujemy to samo
	rozumowanie do przedziału dopełniającego, który ma mniej niż $n$ elementów.
	Zatem pewna krawędź $\ell_i\ell_{i+n}$ łączy obie składowe drzewa bez $e$,
	więc również usunięcie starej krawędzi nie rozspaja grafu.
\end{solutionbox}

\begin{zadanie}{średnie} % Polska 1997; graf ukryty w geometrii
	Na okręgu jednostkowym leży $n$ punktów. Udowodnij, że co najwyżej
	$n^2/3$ par tych punktów jest odległych o więcej niż $\sqrt2$.
\end{zadanie}

\begin{solutionbox}
	Połączmy dwa punkty krawędzią, gdy ich odległość jest większa niż $\sqrt2$.
	Graf nie zawiera $K_4$. Istotnie, ustawmy cztery punkty w kolejności na
	okręgu. Gdyby każda para była połączona, to także każda z czterech par
	kolejnych punktów miałaby cięciwę dłuższą niż $\sqrt2$, a więc odpowiadający
	jej mniejszy kąt środkowy byłby większy niż $90^\circ$. Suma czterech luk
	kątowych byłaby wtedy większa niż $360^\circ$, sprzeczność.

	Pozostaje użyć następującego szczególnego przypadku twierdzenia Turána:
	graf bez $K_4$ na $n$ wierzchołkach ma co najwyżej $n^2/3$ krawędzi.
	Udowodnimy go indukcyjnie. Jeśli graf nie ma trójkąta i ma $m$ krawędzi,
	to dla każdej krawędzi $uv$ zachodzi $d(u)+d(v)\leqslant n$. Sumując po
	krawędziach i używając nierówności Cauchy'ego, otrzymujemy
	\[
		\frac{4m^2}{n}\leqslant\sum_v d(v)^2
		=\sum_{uv\in E(G)}(d(u)+d(v))\leqslant nm,
	\]
	czyli $m\leqslant n^2/4$. Jeśli zaś istnieje trójkąt, to każdy z pozostałych
	wierzchołków ma w nim co najwyżej dwóch sąsiadów. Po usunięciu trzech
	wierzchołków trójkąta i zastosowaniu indukcji liczba krawędzi jest nie
	większa niż
	\[
		\frac{(n-3)^2}{3}+3+2(n-3)=\frac{n^2}{3}.
	\]
	Żądane pary są krawędziami skonstruowanego grafu, więc teza wynika z tego
	oszacowania.
\end{solutionbox}

\begin{zadanie}{trudne} % IMO Shortlist 2004 C3
	Na grafie skończonym wolno wykonać następujący ruch: wybrać cykl długości
	$4$ i usunąć jedną z jego krawędzi. Dla $n\geqslant4$ wyznacz najmniejszą
	liczbę krawędzi, z którą można zakończyć, startując z grafu pełnego $K_n$.
\end{zadanie}

\begin{solutionbox}
	Każda usuwana krawędź leży na cyklu, więc jej usunięcie nie rozspaja grafu.
	Ponadto graf nigdy nie staje się dwudzielny. Istotnie, przypuśćmy, że po
	pierwszym takim ruchu graf jest dwudzielny, a usunęliśmy krawędź $AB$ z
	cyklu $ABCD$. W pozostałym grafie istnieje ścieżka
	$A-D-C-B$ długości $3$, więc $A$ i $B$ leżą w przeciwnych częściach.
	Dołożenie $AB$ zachowałoby dwudzielność, sprzecznie z wyborem pierwszego
	takiego ruchu.

	Graf końcowy jest więc spójny i niedwudzielny. Spójny graf na $n$
	wierzchołkach z $n-1$ krawędziami jest drzewem, a zatem jest dwudzielny.
	Muszą pozostać co najmniej $n$ krawędzi.

	Pokażemy konstrukcję osiągającą $n$. Oznaczmy wierzchołki przez
	$v_1,\dots,v_n$. Najpierw, dla $3\leqslant i<j<n$, usuwamy $v_iv_j$ z
	czworokąta $v_2v_iv_jv_n$. Następnie dla $3\leqslant i<n$ usuwamy
	$v_2v_i$ z cyklu $v_1v_2v_iv_n$ i $v_iv_n$ z cyklu
	$v_1v_iv_nv_2$. Zostają dokładnie krawędzie
	\[
		v_1v_i\quad(2\leqslant i\leqslant n)
		\qquad\text{oraz}\qquad v_2v_n,
	\]
	czyli $n$ krawędzi.
\end{solutionbox}

\begin{zadanie}{trudne} % IMO Shortlist 2002 C6
	Niech $n$ będzie dodatnią liczbą parzystą. Udowodnij, że istnieje permutacja
	$(x_1,\dots,x_n)$ liczb $1,\dots,n$ taka, że dla każdego $i$ (indeksy
	cykliczne) liczba $x_{i+1}$ należy do zbioru
	\[
		\{2x_i,\ 2x_i-1,\ 2x_i-n,\ 2x_i-n-1\}.
	\]
\end{zadanie}

\begin{solutionbox}
	Napiszmy $n=2m$. Budujemy skierowany multigraf o wierzchołkach
	$1,\dots,m$ i krawędziach oznaczonych liczbami $1,\dots,2m$. Krawędź o
	numerze $j$ wychodzi z wierzchołka $\lceil j/2\rceil$ i wchodzi do
	wierzchołka przystającego do $j$ modulo $m$ (z reprezentantem w
	$\{1,\dots,m\}$).

	Z każdego wierzchołka $r$ wychodzą krawędzie $2r-1,2r$, a wchodzą
	krawędzie $r,r+m$. Stopnie wejściowe i wyjściowe są więc równe $2$.
	Graf jest słabo spójny: indukcyjnie można dojść z $1$ do każdego
	$k\leqslant m$, bo do $k$ prowadzi krawędź wychodząca z
	$\lceil k/2\rceil<k$. Istnieje zatem skierowany obwód Eulera.

	Niech $x_1,\dots,x_n$ będą kolejnymi numerami krawędzi na tym obwodzie.
	Jeśli krawędź $x_i$ wchodzi do $r$, to $x_i=r$ albo $r+m$, a następna
	krawędź ma numer $2r$ albo $2r-1$. Po wyeliminowaniu $r$ dostajemy dokładnie
	cztery możliwości z treści zadania.
\end{solutionbox}

% \begin{zadanie}{trudne} % USA TST 2002
% Niech $S$ będzie zbiorem $2^n+1$ elementów, a
% \[
%   f:\binom{S}{2}\longrightarrow\{0,1,\dots,2^{n-1}-1\}.
% \]
% Załóżmy, że dla każdych trzech różnych $x,y,z\in S$ jedna z liczb
% \[
%   f(\{x,y\}),\quad f(\{y,z\}),\quad f(\{z,x\})
% \]
% jest sumą dwóch pozostałych. Udowodnij, że istnieją różne $a,b,c\in S$
% takie, że wszystkie trzy wartości na parach z $\{a,b,c\}$ są równe zero.
% \end{zadanie}

% \begin{solutionbox}
% Prowadzimy indukcję po $n$. Dla $n=1$ wszystkie wartości należą do
% $\{0\}$, więc teza jest oczywista.

% Budujemy graf na $S$, łącząc $x$ i $y$ wtedy i tylko wtedy, gdy
% $f(\{x,y\})$ jest nieparzyste. W każdym trójkącie grafu liczba krawędzi jest
% parzysta: jeśli jedna z trzech wartości jest sumą pozostałych, to suma
% wszystkich trzech wartości jest parzysta. Każda trójka wierzchołków
% wyznacza więc $0$ albo $2$ krawędzie.

% Graf jest dwudzielny. Rzeczywiście, ustalmy wierzchołek $v$. Jego sąsiedzi
% nie mogą mieć między sobą krawędzi, bo z $v$ powstałyby trzy krawędzie.
% Jego niesąsiedzi także nie mogą mieć między sobą krawędzi, bo powstałaby
% dokładnie jedna krawędź. Sąsiedzi i niesąsiedzi tworzą zatem dwie części.

% Większa część $S'$ ma co najmniej
% \[
%   \left\lceil\frac{2^n+1}{2}\right\rceil=2^{n-1}+1
% \]
% elementów. Między jej elementami wszystkie wartości $f$ są parzyste.
% Wybieramy dowolne $2^{n-1}+1$ elementów i definiujemy
% \[
%   g(\{x,y\})=\frac{f(\{x,y\})}{2}.
% \]
% Funkcja $g$ przyjmuje wartości w zbiorze
% $\{0,\dots,2^{n-2}-1\}$ i spełnia ten sam warunek o trójkach. Z założenia
% indukcyjnego istnieje trójka, na której wszystkie wartości $g$, a więc także
% $f$, są zerowe.
% \end{solutionbox}

\begin{zadanie}{bardzo trudne} % IMO 1992, zadanie 3
	Mamy $9$ punktów i wszystkie $36$ łączących je odcinków. Każdy odcinek
	kolorujemy na czerwono, na niebiesko albo pozostawiamy bez koloru. Wyznacz
	najmniejszą liczbę $m$ taką, że każde pokolorowanie dokładnie $m$ odcinków
	wymusza monochromatyczny trójkąt.
\end{zadanie}

\begin{solutionbox}
	Odpowiedzią jest $m=33$.

	Najpierw pokażemy, że $33$ wystarcza. Pozostają wtedy tylko trzy
	niepokolorowane krawędzie. Wybierając po jednym końcu każdej z nich i
	usuwając wybrane wierzchołki, zostawiamy co najmniej sześć wierzchołków,
	między którymi wszystkie krawędzie są pokolorowane. W każdym
	dwukolorowaniu $K_6$ istnieje monochromatyczny trójkąt: z jednego
	wierzchołka wychodzą co najmniej trzy krawędzie jednego koloru; albo któraś
	krawędź między ich końcami ma ten sam kolor, albo wszystkie trzy mają kolor
	przeciwny.

	Pozostaje konstrukcja z $32$ krawędziami. Weźmy punkt $X$ oraz dwa
	rozłączne czworokąty o zbiorach wierzchołków
	\[
		\{R_1,R_2,R_3,R_4\}
		\qquad\text{i}\qquad
		\{B_1,B_2,B_3,B_4\}.
	\]
	Boki pierwszego
	czworokąta kolorujemy na czerwono, drugiego na niebiesko, a ich cztery
	przekątne zostawiamy bez koloru. Krawędzie z $X$ do $R_i$ kolorujemy na
	niebiesko, a z $X$ do $B_i$ na czerwono. Wreszcie krawędź $R_iB_j$
	kolorujemy na czerwono, gdy $i,j$ mają tę samą parzystość, i na niebiesko w
	przeciwnym razie.

	Sprawdźmy trójkąty. Wewnątrz jednego czworokąta każdy trójkąt używa
	niepokolorowanej przekątnej. Trójkąt zawierający $X$ albo ma dwa boki
	różnych kolorów, albo także używa przekątnej. W trójkącie $R_iR_jB_k$,
	gdy $i,j$ mają tę samą parzystość, krawędź $R_iR_j$ jest przekątną; gdy
	mają różną parzystość, krawędzie $R_iB_k$ i $R_jB_k$ mają różne kolory.
	Analogicznie jest dla trójkątów $B_iB_jR_k$. Żaden trójkąt nie jest więc
	monochromatyczny, a pokolorowaliśmy $36-4=32$ krawędzie.
\end{solutionbox}

\begin{zadanie}{bardzo trudne} % AoPS handout, ćw. 20; IMO Shortlist 2013 C6
W pewnym kraju niektóre pary miast są połączone bezpośrednimi lotami w obie
strony. Z każdego miasta można dotrzeć do każdego innego. Wiadomo, że dla
dowolnego miasta istnieje co najwyżej $100$ miast odległych od niego
dokładnie o $3$ loty. Udowodnij, że od żadnego miasta nie może być więcej
niż $2550$ miast odległych dokładnie o $4$ loty.
\end{zadanie}

\begin{solutionbox}
Niech $S_i(v)=\{w:d(v,w)=i\}$. Przypuśćmy przeciwnie, że dla pewnego $x$
zbiór $D=S_4(x)$ ma co najmniej $2551$ elementów, i oznaczmy $A=S_1(x)$.
Wybierzmy najmniejszy zbiór $A^*\subseteq A$ taki, że
\[
  D\subseteq\bigcup_{a\in A^*}S_3(a),
\]
i połóżmy $m=|A^*|$. Ponieważ $m(101-m)\leqslant50\cdot51=2550$, istnieje
$a\in A^*$, dla którego
\[
  |S_3(a)\cap D|\geqslant102-m.
\]
Istotnie, w przeciwnym razie wszystkie $m$ zbiorów pokrywałyby łącznie co
najwyżej $m(101-m)$ elementów $D$.

Zbiór $S_3(a)$ ma najwyżej $100$ elementów. Zatem zbiór
\[
  T=\{c\in S_3(a):d(x,c)\leqslant3\}
\]
ma co najwyżej $m-2$ elementów. Wyprodukujemy w nim jednak $m-1$ różnych
wierzchołków.

Dla każdego $y\in A^*\setminus\{a\}$ minimalność $A^*$ daje miasto
$d_y\in D$, do którego żadna czterolotowa droga z $x$ nie przechodzi przez
wierzchołek z $A^*\setminus\{y\}$. Ustalmy najkrótszą drogę
\[
  x-y-b_y-c_y-d_y.
\]
Wszystkie wierzchołki $b_y,c_y$ są parami różne: równość dwóch
$b$ albo dwóch $c$ pozwoliłaby dojść do odpowiedniego $d_y$ przez inny
element $A^*$, a $b$ nie może być równy $c$, bo leżą w różnych warstwach
BFS od $x$.

Mamy $d(a,y)\leqslant2$. Z drugiej strony $d(a,d_y)\geqslant3$, a równość
nie zachodzi z wyboru $d_y$, więc $d(a,d_y)>3$. Wzdłuż ścieżki
$y-b_y-c_y-d_y$ odległość od $a$ zmienia się o co najwyżej $1$. Dlatego
jedno z miast $b_y,c_y$ leży w $S_3(a)$, a zarazem w odległości co najwyżej
$3$ od $x$, czyli należy do $T$. Dla różnych $y$ dostajemy różne miasta.
Stąd $|T|\geqslant m-1$, sprzecznie z $|T|\leqslant m-2$.
\end{solutionbox}

\begin{zadanie}{bardzo trudne} % AoPS handout, ćw. 21; IMO Shortlist 2004 C8
Dla skończonego grafu $G$ niech $f(G)$ oznacza liczbę trójkątów, a $g(G)$
liczbę podgrafów $K_4$. Wyznacz najmniejszą stałą $c$, dla której
\[
  g(G)^3\leqslant c\,f(G)^4
\]
dla każdego grafu $G$.
\end{zadanie}

\begin{solutionbox}
Odpowiedzią jest $c=3/32$. Najpierw udowodnimy nierówność pomocniczą. Jeśli
graf $H$ ma $E$ krawędzi i $F$ trójkątów, to
\[
  F^2\leqslant\frac{2E^3}{9}. \tag{1}
\]
Dla wierzchołka $v_i$ oznaczmy przez $v_i'$ liczbę jego sąsiadów, a przez
$e_i$ liczbę krawędzi między tymi sąsiadami. Wówczas
\[
  e_i\leqslant\frac{(v_i')^2}{2},\qquad e_i\leqslant E,
  \qquad\text{więc}\qquad e_i\leqslant v_i'\sqrt{E/2}.
\]
Ponieważ $\sum_i v_i'=2E$ i $\sum_i e_i=3F$, po zsumowaniu i podniesieniu
do kwadratu otrzymujemy (1).

Wróćmy do $G$. Niech $G_i$ będzie grafem indukowanym przez sąsiadów
wierzchołka $v_i$, niech $e_i=|E(G_i)|$ i $t_i=f(G_i)$. Mamy
\[
  \sum_i e_i=3f(G),\qquad \sum_i t_i=4g(G).
\]
Zastosowanie (1) do $G_i$ oraz oczywistego $t_i\leqslant f(G)$ daje
\[
  t_i=t_i^{1/3}t_i^{2/3}
  \leqslant f(G)^{1/3}\left(\frac29\right)^{1/3}e_i.
\]
Sumując po $i$, dostajemy
\[
  4g(G)\leqslant3\left(\frac29\right)^{1/3}f(G)^{4/3},
\]
co po podniesieniu do trzeciej potęgi daje
$g(G)^3\leqslant\frac3{32}f(G)^4$.

Stałej nie można zmniejszyć: dla grafów pełnych
\[
  \frac{g(K_N)^3}{f(K_N)^4}
  =\frac{\binom{N}{4}^3}{\binom{N}{3}^4}
  \longrightarrow\frac3{32}.
\]
\end{solutionbox}
\ifshowsolutions
\newpage
\section{Rozwiązania}
\drukujrozwiazania
\fi

\end{document}
